% Options for packages loaded elsewhere
\PassOptionsToPackage{unicode}{hyperref}
\PassOptionsToPackage{hyphens}{url}
\PassOptionsToPackage{dvipsnames,svgnames*,x11names*}{xcolor}
%
\documentclass[
  a4paper,
  UTF8]{ctexart}
\usepackage{amsmath,amssymb}
\usepackage{lmodern}
\usepackage{iftex}
\ifPDFTeX
  \usepackage[T1]{fontenc}
  \usepackage[utf8]{inputenc}
  \usepackage{textcomp} % provide euro and other symbols
\else % if luatex or xetex
  \usepackage{unicode-math}
  \defaultfontfeatures{Scale=MatchLowercase}
  \defaultfontfeatures[\rmfamily]{Ligatures=TeX,Scale=1}
\fi
% Use upquote if available, for straight quotes in verbatim environments
\IfFileExists{upquote.sty}{\usepackage{upquote}}{}
\IfFileExists{microtype.sty}{% use microtype if available
  \usepackage[]{microtype}
  \UseMicrotypeSet[protrusion]{basicmath} % disable protrusion for tt fonts
}{}
\makeatletter
\@ifundefined{KOMAClassName}{% if non-KOMA class
  \IfFileExists{parskip.sty}{%
    \usepackage{parskip}
  }{% else
    \setlength{\parindent}{0pt}
    \setlength{\parskip}{6pt plus 2pt minus 1pt}}
}{% if KOMA class
  \KOMAoptions{parskip=half}}
\makeatother
\usepackage{xcolor}
\IfFileExists{xurl.sty}{\usepackage{xurl}}{} % add URL line breaks if available
\IfFileExists{bookmark.sty}{\usepackage{bookmark}}{\usepackage{hyperref}}
\hypersetup{
  pdftitle={需求文档},
  pdfauthor={DBMS 控制台项目},
  colorlinks=true,
  linkcolor={blue},
  filecolor={Maroon},
  citecolor={Blue},
  urlcolor={blue},
  pdfcreator={LaTeX via pandoc}}
\urlstyle{same} % disable monospaced font for URLs
\usepackage[margin=2.5cm]{geometry}
\usepackage{color}
\usepackage{fancyvrb}
\newcommand{\VerbBar}{|}
\newcommand{\VERB}{\Verb[commandchars=\\\{\}]}
\DefineVerbatimEnvironment{Highlighting}{Verbatim}{commandchars=\\\{\}}
% Add ',fontsize=\small' for more characters per line
\newenvironment{Shaded}{}{}
\newcommand{\AlertTok}[1]{\textcolor[rgb]{1.00,0.00,0.00}{\textbf{#1}}}
\newcommand{\AnnotationTok}[1]{\textcolor[rgb]{0.38,0.63,0.69}{\textbf{\textit{#1}}}}
\newcommand{\AttributeTok}[1]{\textcolor[rgb]{0.49,0.56,0.16}{#1}}
\newcommand{\BaseNTok}[1]{\textcolor[rgb]{0.25,0.63,0.44}{#1}}
\newcommand{\BuiltInTok}[1]{#1}
\newcommand{\CharTok}[1]{\textcolor[rgb]{0.25,0.44,0.63}{#1}}
\newcommand{\CommentTok}[1]{\textcolor[rgb]{0.38,0.63,0.69}{\textit{#1}}}
\newcommand{\CommentVarTok}[1]{\textcolor[rgb]{0.38,0.63,0.69}{\textbf{\textit{#1}}}}
\newcommand{\ConstantTok}[1]{\textcolor[rgb]{0.53,0.00,0.00}{#1}}
\newcommand{\ControlFlowTok}[1]{\textcolor[rgb]{0.00,0.44,0.13}{\textbf{#1}}}
\newcommand{\DataTypeTok}[1]{\textcolor[rgb]{0.56,0.13,0.00}{#1}}
\newcommand{\DecValTok}[1]{\textcolor[rgb]{0.25,0.63,0.44}{#1}}
\newcommand{\DocumentationTok}[1]{\textcolor[rgb]{0.73,0.13,0.13}{\textit{#1}}}
\newcommand{\ErrorTok}[1]{\textcolor[rgb]{1.00,0.00,0.00}{\textbf{#1}}}
\newcommand{\ExtensionTok}[1]{#1}
\newcommand{\FloatTok}[1]{\textcolor[rgb]{0.25,0.63,0.44}{#1}}
\newcommand{\FunctionTok}[1]{\textcolor[rgb]{0.02,0.16,0.49}{#1}}
\newcommand{\ImportTok}[1]{#1}
\newcommand{\InformationTok}[1]{\textcolor[rgb]{0.38,0.63,0.69}{\textbf{\textit{#1}}}}
\newcommand{\KeywordTok}[1]{\textcolor[rgb]{0.00,0.44,0.13}{\textbf{#1}}}
\newcommand{\NormalTok}[1]{#1}
\newcommand{\OperatorTok}[1]{\textcolor[rgb]{0.40,0.40,0.40}{#1}}
\newcommand{\OtherTok}[1]{\textcolor[rgb]{0.00,0.44,0.13}{#1}}
\newcommand{\PreprocessorTok}[1]{\textcolor[rgb]{0.74,0.48,0.00}{#1}}
\newcommand{\RegionMarkerTok}[1]{#1}
\newcommand{\SpecialCharTok}[1]{\textcolor[rgb]{0.25,0.44,0.63}{#1}}
\newcommand{\SpecialStringTok}[1]{\textcolor[rgb]{0.73,0.40,0.53}{#1}}
\newcommand{\StringTok}[1]{\textcolor[rgb]{0.25,0.44,0.63}{#1}}
\newcommand{\VariableTok}[1]{\textcolor[rgb]{0.10,0.09,0.49}{#1}}
\newcommand{\VerbatimStringTok}[1]{\textcolor[rgb]{0.25,0.44,0.63}{#1}}
\newcommand{\WarningTok}[1]{\textcolor[rgb]{0.38,0.63,0.69}{\textbf{\textit{#1}}}}
\usepackage{longtable,booktabs,array}
\usepackage{calc} % for calculating minipage widths
% Correct order of tables after \paragraph or \subparagraph
\usepackage{etoolbox}
\makeatletter
\patchcmd\longtable{\par}{\if@noskipsec\mbox{}\fi\par}{}{}
\makeatother
% Allow footnotes in longtable head/foot
\IfFileExists{footnotehyper.sty}{\usepackage{footnotehyper}}{\usepackage{footnote}}
\makesavenoteenv{longtable}
\setlength{\emergencystretch}{3em} % prevent overfull lines
\providecommand{\tightlist}{%
  \setlength{\itemsep}{0pt}\setlength{\parskip}{0pt}}
\setcounter{secnumdepth}{5}
\ifLuaTeX
  \usepackage{selnolig}  % disable illegal ligatures
\fi

\title{需求文档}
\author{DBMS 控制台项目}
\date{}

\begin{document}
\maketitle

\renewcommand*\contentsname{目录}
{
\hypersetup{linkcolor=}
\setcounter{tocdepth}{3}
\tableofcontents
}
\hypertarget{ux9700ux6c42ux6587ux6863}{%
\section{需求文档}\label{ux9700ux6c42ux6587ux6863}}

\hypertarget{ux6587ux6863ux76eeux7684}{%
\subsection{1. 文档目的}\label{ux6587ux6863ux76eeux7684}}

本文档从产品视角描述 DBMS
控制台项目的建设目标、用户价值、业务边界、核心功能、非功能要求和版本规划。它用于统一项目参与者对产品方向的理解，并作为《需求规格说明书》的上层依据。

\hypertarget{ux9879ux76eeux80ccux666f}{%
\subsection{2. 项目背景}\label{ux9879ux76eeux80ccux666f}}

当前项目已具备基础登录注册、个性化工作台、数据库管理、表结构管理和表数据
CRUD 能力，后续目标是在本地 MySQL 之上继续完善 GUI
查询、统计图表和更完整的审计体验。系统不重新实现数据库内核，而是在本地
MySQL 之上提供安全、易用、可视化、可个性化的数据库管理体验。

\hypertarget{ux4ea7ux54c1ux5b9aux4f4d}{%
\subsection{3. 产品定位}\label{ux4ea7ux54c1ux5b9aux4f4d}}

系统定位为：

\begin{Shaded}
\begin{Highlighting}[]
\NormalTok{面向个人用户的本地 MySQL 可视化管理与数据工作台}
\end{Highlighting}
\end{Shaded}

它同时具备两类属性：

\begin{itemize}
\tightlist
\item
  DBMS
  控制台：管理用户自己的数据库、表结构、索引、约束、视图、表数据和查询。
\item
  个人工作台：支持头像、偏好、背景、主题、统计卡片和面板布局。
\end{itemize}

\hypertarget{ux7528ux6237ux4e0eux4f7fux7528ux573aux666f}{%
\subsection{4.
用户与使用场景}\label{ux7528ux6237ux4e0eux4f7fux7528ux573aux666f}}

\hypertarget{ux76eeux6807ux7528ux6237}{%
\subsubsection{4.1 目标用户}\label{ux76eeux6807ux7528ux6237}}

第一阶段目标用户为单机环境下的个人学习用户或课程项目使用者。

\hypertarget{ux5178ux578bux573aux666f}{%
\subsubsection{4.2 典型场景}\label{ux5178ux578bux573aux666f}}

\begin{longtable}[]{@{}ll@{}}
\toprule
场景 & 用户目标 \\
\midrule
\endhead
登录系统 & 进入个人数据库工作台 \\
创建数据库 & 为课程、实验或业务数据建立独立数据库 \\
设计表结构 & 用 GUI 创建字段、主键、外键、索引和约束 \\
浏览表数据 & 分页查看、排序、筛选表数据 \\
修改表数据 & 在表格中编辑草稿，点击更新后提交数据库 \\
构造查询 & 通过 GUI 完成 JOIN、聚合、分组、子查询等 SELECT 查询 \\
查看统计 & 通过图表了解用户、数据库、表的使用情况 \\
个性化工作台 & 设置头像、背景、主题色和面板卡片 \\
\bottomrule
\end{longtable}

\hypertarget{ux4ea7ux54c1ux8303ux56f4}{%
\subsection{5. 产品范围}\label{ux4ea7ux54c1ux8303ux56f4}}

\hypertarget{ux8303ux56f4ux5185}{%
\subsubsection{5.1 范围内}\label{ux8303ux56f4ux5185}}

\begin{itemize}
\tightlist
\item
  用户注册、登录、双 Token 会话设计。
\item
  用户资料、头像、偏好设置。
\item
  自定义背景、主题色和 Dashboard 卡片。
\item
  用户自己的数据库管理。
\item
  表、字段、索引、约束、视图管理。
\item
  表数据分页浏览、新增、修改、删除。
\item
  行内编辑草稿与显式更新提交。
\item
  GUI SELECT 查询构造器，支持 JOIN、聚合、分组、HAVING、子查询。
\item
  事务会话设计。
\item
  用户级、数据库级、表级统计。
\item
  审计日志、traceId、错误追踪。
\item
  容量限制、频率限制、安全防御。
\end{itemize}

\hypertarget{ux8303ux56f4ux5916}{%
\subsubsection{5.2 范围外}\label{ux8303ux56f4ux5916}}

第一阶段暂不实现：

\begin{itemize}
\tightlist
\item
  多用户共享同一个数据库。
\item
  远程 MySQL 实例接入。
\item
  MySQL 用户账号和权限系统管理。
\item
  任意 SQL 编辑器。
\item
  递归查询、窗口函数、CTE、UNION 等高级 SQL 构造。
\item
  云端对象存储。
\end{itemize}

\hypertarget{ux4ea7ux54c1ux7ed3ux6784}{%
\subsection{6. 产品结构}\label{ux4ea7ux54c1ux7ed3ux6784}}

\begin{Shaded}
\begin{Highlighting}[]
\NormalTok{DBMS 控制台}
\NormalTok{+{-} 用户系统}
\NormalTok{|  +{-} 注册登录}
\NormalTok{|  +{-} 双 Token 会话}
\NormalTok{|  +{-} 用户资料}
\NormalTok{|  +{-} 偏好设置}
\NormalTok{|  +{-} 资源管理}
\NormalTok{+{-} 个人工作台}
\NormalTok{|  +{-} 顶部栏}
\NormalTok{|  +{-} 数据库导航树}
\NormalTok{|  +{-} 自定义面板}
\NormalTok{|  +{-} 背景和主题}
\NormalTok{|  +{-} 最近操作}
\NormalTok{+{-} 数据库管理}
\NormalTok{|  +{-} 创建数据库}
\NormalTok{|  +{-} 修改数据库属性}
\NormalTok{|  +{-} 删除数据库}
\NormalTok{|  +{-} 数据库统计}
\NormalTok{+{-} 表结构管理}
\NormalTok{|  +{-} 表}
\NormalTok{|  +{-} 字段}
\NormalTok{|  +{-} 索引}
\NormalTok{|  +{-} 约束}
\NormalTok{|  +{-} 视图}
\NormalTok{+{-} 数据管理}
\NormalTok{|  +{-} 分页浏览}
\NormalTok{|  +{-} 新增数据}
\NormalTok{|  +{-} 行内草稿编辑}
\NormalTok{|  +{-} 显式更新提交}
\NormalTok{|  +{-} 删除数据}
\NormalTok{+{-} 查询构造器}
\NormalTok{|  +{-} 字段选择}
\NormalTok{|  +{-} 条件组}
\NormalTok{|  +{-} JOIN}
\NormalTok{|  +{-} 聚合分组}
\NormalTok{|  +{-} HAVING}
\NormalTok{|  +{-} 子查询}
\NormalTok{+{-} 统计与审计}
\NormalTok{   +{-} 统计图表}
\NormalTok{   +{-} 操作趋势}
\NormalTok{   +{-} 审计日志}
\NormalTok{   +{-} traceId 追踪}
\end{Highlighting}
\end{Shaded}

\hypertarget{ux6838ux5fc3ux4ea7ux54c1ux89c4ux5219}{%
\subsection{7.
核心产品规则}\label{ux6838ux5fc3ux4ea7ux54c1ux89c4ux5219}}

\hypertarget{ux6570ux636eux9694ux79bb}{%
\subsubsection{7.1 数据隔离}\label{ux6570ux636eux9694ux79bb}}

每个用户只能操作自己的数据库。前端只传
\texttt{databaseId}，后端根据当前登录用户和系统库映射关系解析真实物理库名。

\hypertarget{ux6570ux636eux4feeux6539ux786eux8ba4}{%
\subsubsection{7.2
数据修改确认}\label{ux6570ux636eux4feeux6539ux786eux8ba4}}

用户在网页表格中直接修改表项时，修改只进入本地草稿态。只有点击行旁``更新/提交''按钮后，系统才向后端提交修改并写入数据库。

\hypertarget{ux67e5ux8be2ux5b89ux5168}{%
\subsubsection{7.3 查询安全}\label{ux67e5ux8be2ux5b89ux5168}}

灵活查询通过结构化 AST 实现，不开放任意 SQL 字符串。查询 AST 支持多表
JOIN、聚合、分组、HAVING、派生表子查询、\texttt{IN}
子查询、\texttt{EXISTS} 子查询和标量子查询。

\hypertarget{ux8d44ux6e90ux5b58ux50a8}{%
\subsubsection{7.4 资源存储}\label{ux8d44ux6e90ux5b58ux50a8}}

头像、背景图等用户资源以文件形式保存，系统库仅保存资源元信息和归属关系，不建议将图片二进制直接保存到数据库表中。

\hypertarget{ux5bb9ux91cfux9650ux5236}{%
\subsubsection{7.5 容量限制}\label{ux5bb9ux91cfux9650ux5236}}

系统必须限制数据库数量、表数量、索引数量、查询复杂度、批量写入规模、文件大小和用户资源总量，避免本地
MySQL 和前端页面被误操作拖垮。

\hypertarget{ux7248ux672cux89c4ux5212}{%
\subsection{8. 版本规划}\label{ux7248ux672cux89c4ux5212}}

\begin{longtable}[]{@{}
  >{\raggedright\arraybackslash}p{(\columnwidth - 4\tabcolsep) * \real{0.33}}
  >{\raggedright\arraybackslash}p{(\columnwidth - 4\tabcolsep) * \real{0.33}}
  >{\raggedright\arraybackslash}p{(\columnwidth - 4\tabcolsep) * \real{0.33}}@{}}
\toprule
阶段 & 目标 & 主要交付 \\
\midrule
\endhead
V1 & 认证与工作台地基 & 双 Token、用户信息、Home 骨架、偏好基础 \\
V2 & 个性化工作台 & 主题、背景、响应式布局、首页体验 \\
V3 & 数据库管理 & 数据库列表、创建、修改、删除、容量限制 \\
V4 & 表结构管理 & 表列表、创建表、结构管理入口、表数据预览、删除表 \\
V5 & 数据 CRUD &
分页预览、字段筛选、排序、列显示控制、单行新增、行内编辑、显式提交、单行删除、批量新增、批量修改、批量删除；无主键表浏览和新增、编辑删除只读保护 \\
V6 & GUI 查询 & JOIN、聚合、分组、HAVING、子查询 \\
V7 & 统计与 DIY & 图表、审计、Dashboard 卡片、布局偏好 \\
V5+ & 前期遗留补齐 &
已补齐用户资料资源上传、字段增删改、索引管理、外键接口、字符串有限取值/枚举式检查约束、批量写入接口原子性 \\
\bottomrule
\end{longtable}

\hypertarget{ux6210ux529fux6807ux51c6}{%
\subsection{9. 成功标准}\label{ux6210ux529fux6807ux51c6}}

项目达到以下标准时，认为核心目标落地：

\begin{itemize}
\tightlist
\item
  用户能登录并进入自己的个性化工作台。
\item
  用户能创建和管理自己的 MySQL 数据库。
\item
  用户能通过 GUI 创建和修改表结构。
\item
  用户能安全地浏览和修改表数据。
\item
  用户能通过 GUI 构造常用复杂 SELECT 查询。
\item
  用户能看到用户级、数据库级、表级统计图表。
\item
  系统具备认证、授权、SQL 防护、容量限制、审计日志和错误追踪。
\end{itemize}

\hypertarget{ux5173ux8054ux6587ux6863}{%
\subsection{10. 关联文档}\label{ux5173ux8054ux6587ux6863}}

\begin{longtable}[]{@{}ll@{}}
\toprule
文档 & 说明 \\
\midrule
\endhead
\texttt{需求规格说明书} & 详细需求和验收标准 \\
\texttt{概要设计说明书} & 系统总体设计 \\
\texttt{详细设计说明书} & 模块、流程、伪代码和落地设计 \\
\texttt{接口文档} & API 契约 \\
\bottomrule
\end{longtable}

\end{document}
